\section{Predictive uncertainty}
\label{sec:pi}

Capacity planning needs a distribution, not a point. Reporting a single
average percentage error conflates two sources of uncertainty --- estimation
error in the fitted coefficients, which grows with the extrapolation distance,
and irreducible run-to-run dispersion, which does not. We separate them.

\subsection{Lognormal predictive distribution}

Let $\vartheta=(\log\Tzero,\gamma,\beta)^{\!\top}$ be estimated
from~\eqref{eq:ols} with covariance estimate $\widehat{\Sigma}$, and let
$u_{t}=(1,\log\cc_{t},-\log\n_{t})^{\!\top}$ be the design vector of a target
configuration. In log-space the predictor is linear, $\hat y_{t}=u_{t}^{\!\top}
\hat\vartheta$, so no linearization is required for the mean and
\begin{equation}
\Var(\hat y_{t})=u_{t}^{\!\top}\widehat{\Sigma}\,u_{t},
\qquad
s^{2}(\n_{t},\cc_{t})=u_{t}^{\!\top}\widehat{\Sigma}\,u_{t}+\hat\sigma^{2}.
\label{eq:predvar}
\end{equation}
The first term is estimation variance; it grows quadratically in
$\log\n_{t}$ and therefore penalizes long extrapolations. The second is the
residual variance of a single future run.

\begin{State}[Predictive interval]
\label{st:pi}
If the log-space residuals are approximately Gaussian, a two-sided predictive
interval of nominal level $1-\varepsilon$ for the throughput at
$(\n_{t},\cc_{t})$ is
\begin{equation}
\PI_{1-\varepsilon}
=\Bigl[\exp\bigl(\hat y_{t}-\zq{1-\varepsilon/2}\,s\bigr),\ 
       \exp\bigl(\hat y_{t}+\zq{1-\varepsilon/2}\,s\bigr)\Bigr],
\label{eq:pi}
\end{equation}
with $s=s(\n_{t},\cc_{t})$ from~\eqref{eq:predvar}. The corresponding one-sided
lower bound at level $1-\epsp$ is
$\underline{\TPS}_{1-\epsp}=\exp(\hat y_{t}-\zq{1-\epsp}s)$.
\end{State}

The one-sided form is the quantity entering the sizing problem of
Section~\ref{sec:sizing}: a capacity requirement is a lower-tail statement.

\subsection{Residual bootstrap}

Gaussianity of log-residuals is an assumption, not a fact. We therefore also
compute a distribution-free interval by the residual bootstrap: resample the
centred log-residuals $\{r_{i}-\bar r\}$ with replacement, refit
\eqref{eq:ols} on the resampled responses to obtain $\vartheta^{(b)}$,
and form
\begin{equation}
y_{t}^{(b)}=u_{t}^{\!\top}\vartheta^{(b)}+r^{(b)},\qquad b=1,\dots,B,
\label{eq:bootstrap}
\end{equation}
where $r^{(b)}$ is an independently drawn residual. The empirical quantiles of
$\exp(y_{t}^{(b)})$ give the interval. We report both constructions; agreement
between them is evidence that the Gaussian approximation is adequate.

\subsection{Validation by coverage, not by average error}

The correctness criterion for an interval is its \emph{coverage}: the fraction
of held-out observations that fall inside a nominal $1-\varepsilon$ interval
should be close to $1-\varepsilon$. We therefore validate as follows.

\begin{enumerate*}
\item Calibrate Algorithm~\ref{alg:calibration} using only validator counts
$\n\le\n_{\mathrm{cal}}$.
\item Predict at held-out counts $\n_{t}>\n_{\mathrm{cal}}$ and construct
$\PI_{0.90}$ and $\PI_{0.95}$.
\item Report the empirical coverage over all held-out runs.
\end{enumerate*}

Coverage substantially below nominal indicates that the extrapolation is
optimistic; coverage substantially above nominal indicates unnecessarily wide
intervals and hence overprovisioning. Results appear in
Section~\ref{sec:results}.
